\documentclass[11pt]{article}
\usepackage[margin=1in]{geometry}
\usepackage{amsmath,amssymb}
\usepackage{booktabs}
\usepackage[dvipsnames]{xcolor}
\usepackage{hyperref}
\hypersetup{colorlinks=true,linkcolor=blue!50!black,urlcolor=blue!50!black,citecolor=blue!50!black}

\title{Finite-Height Rigidity of the Riemann Zeta Zeros:\\
A Preregistered Numerical Test Against Random-Matrix Statistics}
\author{David Lindgreen}
\date{October 9, 2026 \\[2pt] \small Version 0.1.0 \quad Preprint, not peer reviewed}

\begin{document}
\maketitle

\begin{abstract}
We test whether two blocks of 2{,}000 consecutive nontrivial zeros of the Riemann zeta function, the
lowest and those near the 100{,}000th, follow the statistics of the Gaussian Unitary Ensemble (GUE).
Hypotheses and thresholds were fixed in a protocol committed before any result was generated. The zeros are
5 to 12 times closer to GUE than to independent random points in nearest-neighbour spacing and pair
correlation, with strong level repulsion. However, the spacing variance is $0.148$ (low block) and
$0.162$ (high block), below the $95\%$ sampling band $[0.169, 0.195]$ of 150 simulated GUE samples of the same
size and below the exact GUE value $0.18000$: the zeros are measurably \emph{more rigid} than finite-size GUE,
and the discrepancy shrinks with height. This direction is consistent with published finite-height corrections
to the Montgomery--Odlyzko law; that model is not tested here. The study is numerical and makes no claim about the
Riemann Hypothesis.
\end{abstract}

\section{Question and context}
Montgomery~\cite{montgomery1973} showed, assuming the Riemann Hypothesis, that the pair correlation of
zeta zeros agrees with the GUE form $R_2(r)=1-(\sin\pi r/\pi r)^2$ for test functions with restricted Fourier
support, and Odlyzko~\cite{odlyzko1987} found close numerical agreement at large height. The agreement is an
asymptotic statement. We ask a bounded finite question: for two disjoint blocks of 2{,}000 zeros, how far are
the empirical nearest-neighbour spacing law and pair correlation from the GUE predictions, compared with
(a) a Poisson control, (b) the sampling scatter of genuine GUE data of the same size, and (c) each other.

\section{Data and unfolding}
Heights $\gamma_n$ of the zeros $n=1,\dots,2000$ (heights $14.13$ to $2515.29$) and $n=100000,\dots,101999$
(heights $74920.83$ to $76257.69$) were computed with \texttt{mpmath} and frozen.
Completeness was verified by Hardy's function $Z(t)$, which changes sign at each simple critical-line zero: the
sign of $Z$ at the midpoint of every gap strictly alternates on all 1{,}999 gaps of both blocks, which a single
missed or duplicated zero always breaks. Zeros are unfolded by $x_n=\theta(\gamma_n)/\pi+1$ with $\theta$ the
Riemann--Siegel theta function~\cite{edwards1974}, so the spacings $s_n=x_{n+1}-x_n$ have mean $1$.

\section{Reference laws and null}
The GUE nearest-neighbour spacing law is computed exactly from the sine-kernel Fredholm determinant
$E(s)=\det(I-K_s)$ by Gauss--Legendre quadrature~\cite{bornemann2010}, with $F(s)=1+E'(s)$; its mean is $1.0000000$
and variance $0.17999$. The Poisson control has $F(s)=1-e^{-s}$ and $R_2\equiv1$. Because one zeta block is a single
realization, every statistic is also computed on 150 simulated GUE samples (seed $20260920$; central,
semicircle-unfolded eigenvalues of independent $800\times800$ matrices pooled to 1{,}999 spacings), giving a calibrated
sampling band. Kolmogorov--Smirnov (KS) values are used only as distances, never $p$-values, since zeta spacings are
not independent.

\section{Results}
\begin{table}[h]
\centering
\begin{tabular}{lrrrr}
\toprule
Statistic & Low block & High block & Exact GUE & GUE band (95\%)\\
\midrule
Spacing KS distance to GUE & 0.0426 & 0.0255 & --- & $\le 0.0243$\\
Pair-correlation RMS to GUE & 0.0681 & 0.0606 & --- & $\le 0.0703$\\
Spacing variance & 0.1479 & 0.1618 & 0.1800 & $0.1686$--$0.1946$\\
Share of spacings $<0.25$ & 0.0080 & 0.0110 & 0.0163 & $0.0115$--$0.0225$\\
\bottomrule
\end{tabular}
\caption{Measured statistics against exact GUE and the simulated GUE sampling band. Poisson distances are
$0.3161$ and $0.3052$ (KS) and $0.3608$ and $0.3489$ (pair-correlation RMS) for the low and high blocks.}
\end{table}

\paragraph{Preregistered hypotheses.}
\textbf{H1} (GUE closer than Poisson on both distances in both blocks): confirmed.
\textbf{H2} (fewer than $7.37\%$ of spacings below $0.25$): confirmed ($0.80\%$, $1.10\%$).
\textbf{H3} (each distance inside the GUE sampling band): \emph{not confirmed}; the pair-correlation RMS passes in
both blocks, while the spacing KS distance exceeds the band (clearly in the low block, marginally in the high block)
and the spacing variance lies below the band in both.
\textbf{H4} (agreement improves with height on both distances): confirmed.

\section{Interpretation and limits}
The excess rigidity and its decrease with height are the direction expected if, at finite height, the zeros behave
like a unitary ensemble of finite effective dimension growing logarithmically with height, as developed by
Bogomolny, Bohigas, Leboeuf and Monastra~\cite{bogomolny2006}. We did not fit or test that model. Further limits:
two blocks only; spacing and pair correlation only; smooth-density unfolding only; the simulated band carries its own
sampling error of a few percent, so the high block's marginal KS exceedance ($0.0255$ against $0.0243$) should
be read cautiously; and each block is one realization, so the height trend establishes direction, not magnitude.
No claim is made about the Riemann Hypothesis.

\section{Proposed follow-up (not a result)}
Preregister a v0.2 that replaces the infinite-$N$ GUE reference by a circular unitary ensemble of effective dimension
set from the published formula without fitting, and test whether both blocks then fall inside that model's sampling
band.

\section*{Reproducibility and AI assistance}
Code, frozen zeros, protocol, registry and tests:\\
\url{https://github.com/lindgreendavid/zeta-zeros-lab}\\
Interactive laboratory:\\
\url{https://zeta-zeros-lab-interactive.lindgreendavid.workers.dev}\\
The protocol commit precedes the
registry commit in the repository history. The analysis code, tests and drafts of this text were produced with
substantial assistance from an AI system (Claude, Anthropic); the author is responsible for the content.

\begin{thebibliography}{9}
\bibitem{montgomery1973} H.~L. Montgomery, The pair correlation of zeros of the zeta function,
\emph{Proc. Sympos. Pure Math.} \textbf{24} (1973), 181--193.
\bibitem{odlyzko1987} A.~M. Odlyzko, On the distribution of spacings between zeros of the zeta function,
\emph{Math. Comp.} \textbf{48} (1987), 273--308.
\bibitem{bornemann2010} F.~Bornemann, On the numerical evaluation of distributions in random matrix theory: a review,
\emph{Markov Process. Related Fields} \textbf{16} (2010), 803--866.
\bibitem{bogomolny2006} E.~B. Bogomolny, O.~Bohigas, P.~Leboeuf, A.~G. Monastra, On the spacing distribution of the
Riemann zeros: corrections to the asymptotic result, \emph{J. Phys. A} \textbf{39} (2006), 10743--10754.
\bibitem{edwards1974} H.~M. Edwards, \emph{Riemann's Zeta Function}, Academic Press, 1974.
\end{thebibliography}
\end{document}
